\documentclass{ximera}

\title{Vraagstukken potentiële energie, potentiaal en spanning: basisoefeningen}
\author{Daan Lipkens}

\begin{document}
\begin{abstract}
\end{abstract}
\maketitle

%———————————————————————————————————————
% BASIS OEFENINGEN POTENTIELE ENERGIE, POTENTIAAL EN SPANNING
%———————————————————————————————————————
\section{Basisoefeningen potentiële energie, potentiaal en spanning}

%———————————————————————————————————————
% BASIS OEFENING 1: Epot, positieve lading
\begin{exercise}
    Een lading $Q = +4{,}0 \cdot 10^{-6} \, \mathrm{C}$ bevindt zich op een afstand van $0{,}25 \, \mathrm{m}$ van de negatieve plaat, in een homogeen elektrisch veld met grootte $E = 300 \, \mathrm{N/C}$. Bereken de potentiële elektrische energie van de lading.

    \[ E_{\text{pot}} = \answer[tolerance=0.01,onlineshowanswerbutton]{3.0} \cdot 10^{-4} \, \mathrm{J} \]

    \begin{solution}\nl
        \underline{Gegevens:}
        \begin{align*}
            & Q = +4{,}0 \cdot 10^{-6} \, \mathrm{C} & & E = 300 \, \mathrm{N/C} & & d = 0{,}25 \, \mathrm{m}
        \end{align*}

        \underline{Gevraagd:} $E_{\text{pot}}$

        \underline{Oplossing:}
        \begin{align*}
            E_{\text{pot}} &= Q \cdot E \cdot d \\
            &= 4{,}0 \cdot 10^{-6} \, \mathrm{C} \cdot 300 \, \mathrm{N/C} \cdot 0{,}25 \, \mathrm{m} \\
            &= 3{,}0 \cdot 10^{-4} \, \mathrm{J}
        \end{align*}
    \end{solution}
\end{exercise}
%———————————————————————————————————————

%———————————————————————————————————————
% BASIS OEFENING 2: Epot, negatieve lading
\begin{exercise}
    Een lading $Q = -6{,}0 \cdot 10^{-6} \, \mathrm{C}$ bevindt zich op een afstand van $0{,}25 \, \mathrm{m}$ van de negatieve plaat, in hetzelfde homogeen elektrisch veld als de vorige oefening ($E = 300 \, \mathrm{N/C}$). Bereken de potentiële elektrische energie van de lading.

    \[ E_{\text{pot}} = \answer[tolerance=0.01,onlineshowanswerbutton]{-4.5} \cdot 10^{-4} \, \mathrm{J} \]

    \begin{hint}
        Vul het teken van $Q$ mee in de formule; in tegenstelling tot bij de wet van Coulomb neem je hier \textbf{geen} absolute waarde.
    \end{hint}

    \begin{solution}\nl
        \underline{Gegevens:}
        \begin{align*}
            & Q = -6{,}0 \cdot 10^{-6} \, \mathrm{C} & & E = 300 \, \mathrm{N/C} & & d = 0{,}25 \, \mathrm{m}
        \end{align*}

        \underline{Gevraagd:} $E_{\text{pot}}$

        \underline{Oplossing:}
        \begin{align*}
            E_{\text{pot}} &= Q \cdot E \cdot d \\
            &= -6{,}0 \cdot 10^{-6} \, \mathrm{C} \cdot 300 \, \mathrm{N/C} \cdot 0{,}25 \, \mathrm{m} \\
            &= -4{,}5 \cdot 10^{-4} \, \mathrm{J}
        \end{align*}
        De potentiële energie is negatief omdat de lading zelf negatief is.
    \end{solution}
\end{exercise}
%———————————————————————————————————————

%———————————————————————————————————————
% BASIS OEFENING 3: V = E*d
\begin{exercise}
    Bereken de potentiaal in een punt op $0{,}12 \, \mathrm{m}$ van de negatieve plaat, in een homogeen elektrisch veld met grootte $E = 400 \, \mathrm{N/C}$.

    \[ V = \answer[tolerance=0.5,onlineshowanswerbutton]{48} \, \mathrm{V} \]

    \begin{solution}\nl
        \underline{Gegevens:}
        \begin{align*}
            & E = 400 \, \mathrm{N/C} & & d = 0{,}12 \, \mathrm{m}
        \end{align*}

        \underline{Gevraagd:} $V$

        \underline{Oplossing:}
        \begin{align*}
            V = E \cdot d = 400 \, \mathrm{N/C} \cdot 0{,}12 \, \mathrm{m} = 48 \, \mathrm{V}
        \end{align*}
    \end{solution}
\end{exercise}
%———————————————————————————————————————

%———————————————————————————————————————
% BASIS OEFENING 4: U = E*l
\begin{exercise}
    Een geleider met lengte $0{,}40 \, \mathrm{m}$ bevat een homogeen elektrisch veld met grootte $150 \, \mathrm{N/C}$. Bereken de spanning over de geleider.

    \[ U = \answer[tolerance=0.5,onlineshowanswerbutton]{60} \, \mathrm{V} \]

    \begin{solution}\nl
        \underline{Gegevens:}
        \begin{align*}
            & E = 150 \, \mathrm{N/C} & & \ell = 0{,}40 \, \mathrm{m}
        \end{align*}

        \underline{Gevraagd:} $U$

        \underline{Oplossing:}
        \begin{align*}
            U = E \cdot \ell = 150 \, \mathrm{N/C} \cdot 0{,}40 \, \mathrm{m} = 60 \, \mathrm{V}
        \end{align*}
    \end{solution}
\end{exercise}
%———————————————————————————————————————

%———————————————————————————————————————
% BASIS OEFENING 5: omgekeerd, gevraagd E uit V en d
\begin{exercise}
    In een punt op $0{,}30 \, \mathrm{m}$ van de negatieve plaat van een homogeen veld, is de potentiaal $90 \, \mathrm{V}$. Bereken de grootte van het elektrisch veld.

    \[ E = \answer[tolerance=1,onlineshowanswerbutton]{300} \, \mathrm{N/C} \]

    \begin{hint}
        Vorm de formule $V = E \cdot d$ om naar $E$.
    \end{hint}

    \begin{solution}\nl
        \underline{Gegevens:}
        \begin{align*}
            & V = 90 \, \mathrm{V} & & d = 0{,}30 \, \mathrm{m}
        \end{align*}

        \underline{Gevraagd:} $E$

        \underline{Oplossing:}
        \begin{align*}
            E = \frac{V}{d} = \frac{90 \, \mathrm{V}}{0{,}30 \, \mathrm{m}} = 300 \, \mathrm{N/C}
        \end{align*}
    \end{solution}
\end{exercise}
%———————————————————————————————————————

%———————————————————————————————————————
% BASIS OEFENING 6: omgekeerd, gevraagd Q uit Epot, E en d
\begin{exercise}
    Een lading heeft in een homogeen elektrisch veld met grootte $E = 200 \, \mathrm{N/C}$, op $0{,}10 \, \mathrm{m}$ van de negatieve plaat, een potentiële elektrische energie van $3{,}0 \cdot 10^{-4} \, \mathrm{J}$. Bereken de grootte van de lading.

    \[ Q = \answer[tolerance=0.5,onlineshowanswerbutton]{15} \cdot 10^{-6} \, \mathrm{C} \]

    \begin{hint}
        Vorm de formule $E_{\text{pot}} = Q \cdot E \cdot d$ om naar $Q$.
    \end{hint}

    \begin{solution}\nl
        \underline{Gegevens:}
        \begin{align*}
            & E_{\text{pot}} = 3{,}0 \cdot 10^{-4} \, \mathrm{J} & & E = 200 \, \mathrm{N/C} & & d = 0{,}10 \, \mathrm{m}
        \end{align*}

        \underline{Gevraagd:} $Q$

        \underline{Oplossing:}
        \begin{align*}
            Q = \frac{E_{\text{pot}}}{E \cdot d} = \frac{3{,}0 \cdot 10^{-4} \, \mathrm{J}}{200 \, \mathrm{N/C} \cdot 0{,}10 \, \mathrm{m}} = 15 \cdot 10^{-6} \, \mathrm{C}
        \end{align*}
    \end{solution}
\end{exercise}
%———————————————————————————————————————

%———————————————————————————————————————
% BASIS OEFENING 7 (moeilijkst): spanning tussen twee punten die geen van beide op de platen liggen
\begin{exercise}
    In een homogeen elektrisch veld met grootte $E = 500 \, \mathrm{N/C}$ bevindt punt $1$ zich op $0{,}08 \, \mathrm{m}$ van de negatieve plaat, en punt $2$ op $0{,}20 \, \mathrm{m}$ van de negatieve plaat. Bereken de spanning $U_{12}$ tussen deze twee punten.

    \[ U_{12} = \answer[tolerance=0.5,onlineshowanswerbutton]{-60} \, \mathrm{V} \]

    \begin{hint}
        Bereken eerst de potentiaal in elk van de twee punten afzonderlijk, en gebruik daarna de definitie $U_{12} = V(1) - V(2)$.
    \end{hint}

    \begin{solution}\nl
        \underline{Gegevens:}
        \begin{align*}
            & E = 500 \, \mathrm{N/C} & & d_{1} = 0{,}08 \, \mathrm{m} & & d_{2} = 0{,}20 \, \mathrm{m}
        \end{align*}

        \underline{Gevraagd:} $U_{12}$

        \underline{Oplossing:}\\
        We berekenen eerst de potentiaal in elk punt:
        \begin{align*}
            V(1) &= E \cdot d_{1} = 500 \, \mathrm{N/C} \cdot 0{,}08 \, \mathrm{m} = 40 \, \mathrm{V} \\
            V(2) &= E \cdot d_{2} = 500 \, \mathrm{N/C} \cdot 0{,}20 \, \mathrm{m} = 100 \, \mathrm{V}
        \end{align*}
        De spanning tussen de twee punten is dan:
        \begin{align*}
            U_{12} = V(1) - V(2) = 40 \, \mathrm{V} - 100 \, \mathrm{V} = -60 \, \mathrm{V}
        \end{align*}
        Het negatieve teken betekent dat punt $2$ de hoogste potentiaal heeft: de spanningspijl tussen $1$ en $2$ wijst dus eigenlijk naar punt $2$.
    \end{solution}
\end{exercise}
%———————————————————————————————————————

\end{document}